\section{Integral cohomology and determinant comparisons}\label{sec:models}

Fix an odd prime $p$.  The arithmetic comparisons below will take place
over power-series rings with coefficients in $\Z_p$, before any
rationalization.  Three features of these comparisons are essential:
specialization is derived, local conditions retain their torsion, and
the duality maps are retained together with the complexes.  We develop
these features first.  The last part of the section records the
specialization criteria that will turn characterwise comparisons into
integral divisibility.

\subsection{Complexes and finite models}

Put $T=T_pE$ and $V=T\otimes_{\Z_p}\Q_p$.  The Weil pairing identifies
$T$ with its Tate dual $T^*(1)$.  Cohomological Frobenius is arithmetic.
Character weights in a theta sum or an Euler factor may therefore use
the inverse of a cohomological character; each such comparison will
specify its convention.  Character corestriction is the unnormalized
sum of translates.  Its restriction to a field killing the character
is the weighted sum with the character giving the required inverse
equivariance, without division by the degree.

For a number field $F$ and a support $S$ containing the places above $p$
and infinity, write
\[
 G(M)=C^\bullet(G_{F,S},M),\qquad L_v(M)=C^\bullet(F_v,M).
\]
Here $M$ is a finite free coefficient module, initially over a finite
local Artin ring.  A local condition is a map $U_v\longrightarrow L_v$.
With $Q_v=\operatorname{cone}(U_v\to L_v)$, its Selmer complex is
\begin{equation}\label{models:selmer-complex}
 C(U,M)=\operatorname{fib}\left(G(M)\longrightarrow\bigoplus_vQ_v\right),
 \qquad \operatorname{fib}=\operatorname{cone}[-1].
\end{equation}
A full condition is the identity and a strict condition is zero.
Global duality uses modified localization at infinity.  Since $p$ is
odd, real Tate complexes are contractible.  The Euler characteristic
of ordinary global cohomology is the negative of the sum of the
negative eigenspace dimensions at real places and the coefficient
ranks at complex places.

The finite-model construction of \cite[Section 3, ``Finite models and
matrix limits'']{B} has the following integral form.  We
give the construction because the residual specialization used in
Section~\ref{sec:vertical} is not a characteristic-zero formal germ.

\begin{lemma}[Finite models and marked operations]\label{models:finite-models}
Consider continuous-cochain diagrams over finite local Artin rings of
$p$-power characteristic, with finite free coefficient modules and
compatible coefficient lifts.  Assume that residual cohomology
dimensions are bounded uniformly, and vanish outside a uniformly
bounded range for every complex for which a bounded model is required.
Then the diagrams admit bounded finite free models with derived base
change.  Restriction, unramified inflation, local cones, cup products,
invariant maps, conjugation, and specified comparison homotopies can be
retained simultaneously.

Suppose the scalar rings at the stages identify with every fixed Artin
precision of a complete local ring $R$ with finite residue field.  On a
fixed nonprincipal ultrafilter $\mathcal U$, the model matrices have compatible
limits over $R$.  Their marked identities survive arbitrary derived
base change, in particular to residue fields of primes of $R$,
including primes containing $p$.  Fixed-group diagrams whose actions
stabilize at every precision compute the actual cohomology of that
group.  A finite diagram may be enlarged on the same stages while
retaining its previous models and markings.
\end{lemma}

\begin{proof}
Continuous cochains with free coefficients over an Artin ring $A$ are
flat: continuous functions are filtered colimits of finite free modules
indexed by finite clopen partitions.  Tensoring them with a finite
scalar module gives the corresponding continuous cochains.  A flat
module over a local Artin ring is free.  Indeed, lift a basis modulo the
maximal ideal; nilpotence proves surjectivity, and flatness followed by
nilpotence proves that the kernel is zero.

Split the residual complex into its cohomology and contractible disks,
and lift the graded bases.  Write $d_0$ for the disk differential and
$i,q,h$ for the inclusion, projection, and contraction, normalized by
$h^2=hi=qh=0$ and $1-iq=d_0h+hd_0$.  For $d=d_0+\epsilon$, the formulas
\begin{equation}\label{models:contraction}
 \begin{aligned}
 I&=(1+h\epsilon)^{-1}i,&Q&=q(1+\epsilon h)^{-1},\\
 H&=h(1+\epsilon h)^{-1},&
 d_{\rm small}&=q\epsilon(1+h\epsilon)^{-1}i
 \end{aligned}
\end{equation}
are finite series in the nilpotent maximal ideal.  Substitution, using
$d^2=0$, gives $dI=Id_{\rm small}$, $Qd=d_{\rm small}Q$, $QI=1$,
and $dH+Hd=1-IQ$.  The remaining ranks are the residual cohomology
dimensions; hence they have the asserted bounds.  Inserting $I,Q$
transfers maps, and inserting $H$ transfers their composition
homotopies.  Tensoring the contractions transfers the cup maps and
their homotopies.  Cones may be formed from transferred maps and then
contracted.

Functions lift across coefficient quotients, so the contractions can
be chosen compatibly.  At a fixed precision the matrices belong to
finite sets.  Their ultrafilter values are therefore defined, and all
finite matrix identities survive passage to the limit.  Exact
triangles survive by retaining their cone identifications; duality
survives by retaining its maps and contracting the comparison cones.
Because the models are free, ordinary tensor of these models is
derived tensor.  No inversion of $p$ enters this argument.

For a fixed group, compare the limiting model at each precision with
a contraction of its actual cochains.  The two compositions are
homotopic to the identities.  These comparisons preserve the marked
maps.  Passing to compact inverse limits introduces no $\lim^1$ term,
by finiteness and compactness.  When the diagram is enlarged, keep the
old contractions, transfer the additional maps, and decrease the
working precision along the ultrafilter if necessary.  Any finite list
of new identities is then retained at a cofinal sequence of precisions.
\end{proof}

The hypotheses hold for the diagrams used here.  The residual
representations belong to a fixed finite list, and the support has
bounded cardinality.  Pass to a fixed finite extension trivializing
these representations and containing $\mu_p$.  Kummer theory bounds
degree one by the fixed unit and class groups and the bounded number
of allowed valuations.  Inflation--restriction, local and global
duality, and the odd-primary cohomological-dimension bounds give the
remaining bounds.  Thus the number of moving primes matters, but
their ramification degrees need not be bounded.

\subsection{Localization and exact orthogonality}

At a non-$p$ place where $M$ is unramified, the unramified complex is
\[
 U_v^{\rm ur}=[M\xrightarrow{\Frob_v-1}M]\quad\text{in degrees }0,1.
\]
Its quotient in the full local complex is residue cohomology on
$M(-1)$, in degrees $1,2$.  Allowing such a place globally while imposing
$U_v^{\rm ur}$ gives the original global complex, by unramified
inflation.  These are the usual localization and purity comparisons;
they hold for the models with their maps, not just for dimensions of
cohomology.

For completeness, the $G_{F,S}$ comparison can be checked after finite
covers trivializing residual coefficients.  In the universal cover,
degree-one classes die as torsors.  Brauer classes die after adjoining
sufficient $p$-power roots of unity, whose local degrees grow at the
finitely many required places, and Picard classes die after
principalization.  Purity, cohomological dimension, and Brauer
reciprocity give higher vanishing.  Real groups have order prime to
$p$.  This proves the comparison on residual coefficients, and the
finite-model construction lifts it to all the coefficients in use.

Locally, first take invariants under the pro-prime-to-$p$ kernel of
inertia, an exact operation, and then use the procyclic tame
$p$-inertia resolution.  A change of generator $z\mapsto z^b$ acts on
its degree-one term by $(z^b-1)/(z-1)$.  If the conjugating powers tend
to $1$ at the retained precisions, both the powers and these geometric
sums tend to $1$.  This observation fixes the local comparison when
the residue cardinalities of moving primes tend to $1$.

After inversion of $p$, at a place with fixed possibly ramified Tate
action and varying unramified scalar, use residue cochains on $V^{I_v}$.
The singular quotient uses $V_{I_v}(-1)$ in degrees $1,2$.  Inertia has
constant linear algebra before the Frobenius scalar is introduced.
These unramified conditions contain all $H^0$ and inject into field
$H^1$.  Their paired cup is nullhomotopic through residue cochains.
On a field fiber their degree-one dimension is $h^0$, and local
duality and Euler characteristic zero identify their exact
complements.

\begin{lemma}[Exact local complements]\label{models:orthogonality}
Let paired perfect local maps $U_v\to L_v$ and $U'_v\to L'_v$ have
a specified nullhomotopy of their cup restriction.  Suppose the
induced morphism
\[
 L'_v/U'_v\longrightarrow U_v^\vee[-2]
\]
is an equivalence, where $\vee$ denotes derived linear dual before
shifting.  Then their Selmer complexes are perfect duals with shift
$-3$.  The localization boundaries for nested such conditions are
adjoint under this duality.  These assertions hold over the model
rings and after localization or derived specialization.
\end{lemma}

\begin{proof}
Let $G_c$ be the compact-support fiber of global localization.
Poitou--Tate pairs $G_c$ with the unrestricted dual global complex
with shift $-3$, and local Tate duality pairs $L_v,L'_v$ with shift
$-2$.  Dualizing $G_c\to C(U)\to\bigoplus_vU_v$ and using the given
local equivalences identifies the dual fiber with $C(U')$.  The same
triangle identifies the boundary maps as adjoints.

At finite precision these are the cup, invariant, and compact trace
diagrams of local and global duality \cite{Milne,Nekovar}.  The top
scalar local and compact cohomology groups, in degrees two and three,
are the coefficient ring; their traces are compatible with reduction
and corestriction.  Retain these maps and the reciprocity homotopy in
Lemma~\ref{models:finite-models}.  Perfectness is checked on the
residue field and lifted by contracting the comparison cone.  This
also proves compatibility after specialization.  For split local
complexes the required nullhomotopies are checked on the split
cohomology; for unramified conditions they are supplied by the residue
factorization above.
\end{proof}

Conjugate transport may identify a coefficient over an imaginary
quadratic field with its Tate dual.  In that convention complements
are imposed at conjugate places.  Under Shapiro, compact and local
invariants are compatible with corestriction; no factor equal to the
quadratic degree is inserted.  Exact orthogonality over the testing
ring is needed throughout.  Orthogonality only after passage to its
fraction field would not suffice for the determinant comparisons.

\subsection{Detection of classes after residual specialization}

The models also allow classes to be tested on sequences of Galois
elements.  This will let Chebotarev produce new primes whose
first-order effects distinguish specified classes.

\begin{lemma}[Evaluation detection]\label{models:evaluations}
Let $G_R(M)$ be an unrestricted limiting global model as in
Lemma~\ref{models:finite-models}.  After any field specialization,
including $R\to\Frac(\F_p[[\mathbf u]])$, its degree-one cohomology
injects into abstract crossed-homomorphism classes on the product of
the stage absolute Galois groups.  The action is obtained by taking
entrywise coefficient limits before specialization.  Invariants are
detected in the same way.  The assertion holds for a Selmer complex
whenever its local conditions contain all local $H^0$ and inject into
local field $H^1$.
\end{lemma}

\begin{proof}
Keep the inclusions of the finite models into actual cochains.  At
residual precision a bounded list of group elements detects
invariants and detects whether the values of a model cocycle have
the form $g a-a$ for one $a$ in the coefficient module.  The bound
comes from linear equations in the coefficient space and in the
bounded model space.  Mark the resulting evaluation map to
$[M\to M^{\rm list}]$ in degrees zero and one.  It is an isomorphism
on residual $H^0$ and injective on residual $H^1$.  Its cone has zero
residual cohomology in degrees at most zero.  A minimal model of the
cone therefore has zero terms there.  This remains true after every
field base change and proves the required injection.

For every sequence of group elements, and every pair or finite tuple
of sequences, the evaluation matrices have compatible limits.
The cocycle and coboundary equations consequently hold on their
product group.  Conjugation identities hold on classes by the marked
homotopies.  The Selmer assertion follows from the localization
sequence and the hypotheses on the local maps.
\end{proof}

Here are two consequences that we will use.  If a normal joint kernel
$J$ acts trivially on the coefficients, restriction gives additive,
equivariant evaluations on $J$.  It detects classes whenever the
quotient has vanishing $H^1$.  A central scalar whose difference from
$1$ is invertible gives this vanishing directly from the crossed
homomorphism identity.  In characteristic zero another useful
quotient test is a normal subgroup consisting of arbitrary sequences
in a fixed deep $\operatorname{SL}_2(\Z_p)$ subgroup, with its standard
plane action and trivial scalar twists.  To see vanishing of even
abstract $H^1$, use a constant diagonal $(m,m^{-1})$, with integer
$m>1$ sufficiently close to $1$, to kill the cocycle on the diagonal.
Conjugation multiplies the upper-unipotent parameter by $m^2$; the
crossed relation and the $m^2$-power formula then kill the upper
unipotent values, since the resulting coefficient matrix is
invertible.  Use the inverse diagonal for lower unipotents, and
Gaussian factorization for arbitrary sequences.  Such a subgroup
remains available in a non-CM Tate image in the presence of competing
data of bounded derived length: sufficiently iterated derived groups
kill those data and still contain a deep special linear subgroup.

For an absolutely simple coefficient representation with scalar
endomorphisms, evaluations of independent detected classes span the
full tuple of coefficient spaces.  Otherwise their equivariant span,
a submodule of a sum of copies of that simple representation, would
give a scalar relation among the classes.  The same holds jointly for
pairwise nonisomorphic simples.  This statement is valid in the
residual fields; it requires neither characteristic-zero
semisimplicity of the Galois group nor an assertion about polynomial
avoidance on a finite field.  We will use its linear span and
bilinearity there.

Finally, a finite table of the evaluations and other continuous
arithmetic data factors through a finite quotient at each retained
precision.  Chebotarev realizes it by a Frobenius representative once
the prescribed splitting conditions are compatible.  Adding primes
uses unramified inflation with its marked homotopies, so evaluations
already made are preserved.  Cup-primitive equations solved in a
field-specialized model are interpreted on these evaluation cochains
using the same inclusions and homotopies.  Reciprocity in this use is
the transferred arithmetic invariant identity, not a reciprocity
assertion for an abstract product group.  Over a fixed finite extension
of $\Q_p$, fixed denominators can be cleared and equations lifted to
increasing precision.  Such clearing is used only after inversion of
$p$.

\subsection{Determinant volumes and switching lemmas}

For a perfect complex define its inverse determinant by
\[
 \mathcal D(C)=\bigotimes_i(\det C^i)^{(-1)^{i+1}}.
\]
At a DVR $A$, let $d(C,y)$ be the valuation of a generic tensor $y$
relative to a generator of this line.  A torsion cohomology group of
length $b$ in degree $i$ contributes $(-1)^{i+1}b$.  In particular, for
a generically acyclic square complex in degrees $1,2$,
\begin{equation}\label{models:determinant-sign}
 d(C,1)=-v(\det d^1).
\end{equation}
Determinants multiply in exact triangles.  Bases and their signs are
understood up to integral units.

The following statements give the precise interfaces of the algebra
lemmas in \cite[Lemmas 3.5 and 3.6]{B}.  Their hypotheses will be verified
at the primes where the switches are used.

Here is the local diagram used in both statements.  At one place
choose conditions
\[
 U_-\subset U_f,U_s\subset U_+,
 \qquad U_f\cap U_s=U_-,\qquad U_f+U_s=U_+,
\]
whose quotients are split free in degree one:
\[
 U_f/U_-=F[-1],\qquad U_s/U_-=S[-1],\qquad
 U_+/U_-=(F\oplus S)[-1].
\]
The conditions contain all local $H^0$; $U_+$ omits local $H^2$.
Write
$C_-,C_f,C_s,C_+$ for the global complexes with these respective
conditions, keeping all other places fixed.  The localization
maps from $H^1(C_f)$ to $F$ and from $H^1(C_s)$ to $S$ are denoted
by $f$ and $s$.  The dual coefficient has the exactly complementary
diagram, and the mixed pairing between the corresponding finite
and transverse quotients is perfect.  A fiber means derived
tensor with the DVR residue field; a generic line has a primitive
fiber reduction by taking a generator of its saturated kernel in
the two-term free model.  The opposite condition in the line
lemma is the complementary condition on the dual coefficient.

\begin{lemma}[Line switch]\label{models:line-switch}
Over a DVR $A$ with residue field $k$, suppose local finite and transverse coordinates are split
free lines in degree one, with perfect mixed pairing.  The lower
condition contains all $H^0$, the upper condition omits $H^2$, and the
global lower and upper diagrams have exact complementary duality.
Assume the finite and transverse global complexes have amplitude
$[1,2]$ on the generic and residual tests.  Suppose the finite complex
has only a generic degree-one line.  At a fiber with
$\dim H^1(C_f\otimes_A k)=h>1$,
assume nonzero finite evaluation both for a primitive reduction of
that line and for some degree-one class of the opposite condition.
Then switching lowers the degree-one dimension by one, and leaves one
generic degree-one line.  If integral classes $Y,Y'$ satisfy
$s(Y')=f(Y)$ under a unit-volume isomorphism of coordinates, then
\[
 d(C_f,Y)=d(C_s,Y').
\]
Successive switches ending at dimension one prove nonnegativity of
this common valuation.
\end{lemma}

\begin{proof}
The generic finite evaluation is nonzero, so the lower complex is
generically acyclic.  In the fiber, reciprocity against the nonzero
opposite finite evaluation prohibits a transverse contribution.
The two localization triangles therefore lower the dimension by one.
Those same triangles identify the class tensors through the common
lower tensor $1$; the local isomorphism changes its volume only by a
unit.  At dimension one the minimal model is a free line in degree
one, and the integral class has nonnegative coordinate valuation.
\end{proof}

\begin{lemma}[Square switch]\label{models:square-switch}
In the preceding local diagram over a DVR $A$ with residue field
$k$, let $F,S$ be split free planes.  Retain all $H^0$ in the lower
condition and omit $H^2$ from the upper condition.  Assume amplitude
$[1,2]$ for both switched global complexes on their generic and
residual tests, exact self-duality of the lower and upper diagrams,
and perfect mixed pairing.  The pure finite and transverse planes
are isotropic, with their compatible cup nullhomotopies retained.
On a common plane let $e$ be unimodular
alternating and $J$ an integral involution, with cross pairing
$e(x,Jy)$ up to a unit.  Suppose the finite complex has generic rank
one in degrees one and two, with class $Y$ and $f(Y)\ne0$, and an
integral transverse class satisfies $s(Y')=Jf(Y)$ up to a unit.  Then
the transverse complex has the same generic ranks and
\[
 d(C_f,(Y,Y^\vee))=d(C_s,(Y',(Y')^\vee)),
\]
where the dual tensor uses the pairing functional on degree two.
If finite evaluation has rank two at a fiber with
$\dim H^1(C_f\otimes_A k)>1$ and is nonzero
on a primitive reduction of $Y$, the switch lowers its degree-one
dimension by two.  Ending at dimension one gives nonnegative
valuation for the class-functional tensor.
\end{lemma}

\begin{proof}
Adjunction identifies the annihilator of the upper localization
image with that same image, so the image is self-annihilating in the
four coordinates.  The two generic classes form a basis of the upper
degree-one space.  In the dual basis of $H^2(C_-)$, write
$Y^*,(Y')^*$ for the functionals normalized to take value one on
their respective upper basis vectors.  These are dual basis
vectors, distinct from the pairing functionals $Y^\vee,(Y')^\vee$
used in the determinant tensors.  Choose a rational complement
$x$ to $f(Y)$ whose boundary is $(Y')^*$.  The cross-pairing gives
$e(x,J s(Y'))$ a unit value, and $s(Y')=Jf(Y)$ therefore makes
$f(Y)\wedge x$ a unit plane volume.  The complement need not be
integral: if both classes are multiplied by a uniformizer
$\varpi$, the normalized dual $(Y')^*$ and this complement
are divided by $\varpi$.
Interchanging the planes gives
the same volume.  Thus both triangles identify the class-functional
tensors with the same lower tensor, proving the equality.  At a
rank-two fiber the upper image is exactly the finite plane, giving
the dimension drop.  At the terminal fiber the model is two free
lines with zero differential and unimodular pairing.  The tensor is
integral there.  Compatible nullhomotopies on the lower comparisons
are part of exact duality and are retained at all unchanged places.
\end{proof}

\begin{lemma}[Strict versus finite conditions]\label{models:strict-finite}
At a split pair $w,\bar w$, let $C_*$ be strict at $w$ and full at
$\bar w$, and let $C_F$ impose exactly complementary finite
conditions $U_w,U_{\bar w}$.  Keep all other conditions fixed with
conjugate duality.  Suppose these comparisons are perfect over a DVR,
$U_w$ is generically a line in degree one, $C_*$ is generically
acyclic, and $Y\in H^1(C_F)$ has nonzero localization.  If
$\operatorname{loc}_wY=a e$ in a generic basis $e$ of $U_w$, then
\begin{equation}\label{models:strict-finite-volume}
 d(C_F,(Y,Y^\vee))=d(C_*,1)+2\bigl(v(a)+d(U_w,e)\bigr).
\end{equation}
Moreover $C_F$ has exactly one generic line in each of degrees one
and two.  Conversely, these line assertions and the nonzero
localization imply the strict acyclicity.
\end{lemma}

\begin{proof}
Use $C_0$, strict at $w$ and finite at $\bar w$, and the triangles
$C_0\to C_F\to U_w$ and $C_0\to C_*\to L_{\bar w}/U_{\bar w}$.
The latter quotient is $U_w^\vee[-2]$.  Its boundary of $e^*$ pairs
with $Y$ by $a$.  The two triangles therefore compare the tensors
with a factor $a^2$, and their local lattice volumes are dual.  This
gives the formula and the rank assertions; the converse follows by
the same nonzero adjoint boundary.  This is the determinant
calculation of \cite[Lemma 3.4]{B}.
\end{proof}

\subsection{The Kummer lattice and its volume}

At an actual scalar specialization the integral local condition is
\[
 U_v^{\rm Kum}=E(F_v)^{\wedge}_p[-1]
 \longrightarrow C^\bullet(F_v,T).
\]
The completed point module includes its torsion.  Its derived
reduction modulo $p^n$ therefore has the invariant $p^n$-torsion in
degree zero and the finite Kummer image in degree one.  The
Bockstein identifies the former with the integral torsion Kummer
classes.  Finite Kummer orthogonality and derived base change show
that these conditions are exact complements.  More explicitly the
cup restriction is null: morphisms from the derived Kummer tensor
to scalars in degree two are detected in top degree, and finite
Kummer pairings vanish for every $n$.  The resulting local quotient
duality is an equivalence on the residue field and hence integrally.

Compact local duality gives
\[
 0\longrightarrow E(F_v)^{\wedge}_p\longrightarrow H^1(F_v,T)
 \longrightarrow
 \Hom(E^\vee(F_v)^{\wedge}_p,\Z_p)\longrightarrow0.
\]
In particular at a non-$p$ finite place the compact degree-one
cohomology is entirely Kummer.  Over $F_v/\Q_p$, the rational
abelian-variety comparison of Bloch--Kato gives
\begin{equation}\label{models:log-exp-pairing}
 \operatorname{inv}_v(z\cup x)
 =\pm\Tr_{F_v/\Q_p}\bigl(\exp^*z\cdot\log x\bigr).
\end{equation}
For $F_v=\Q_p$ and $\exp^*z=\alpha\omega_E$, this is
$\pm\alpha\log_{\omega_E}x$.  This rational comparison includes bad
reduction; it makes no assertion that exponential coordinates are
integral quotient bases \cite{BlochKato}.

\begin{lemma}[Local Haar factors]\label{models:haar}
For the minimal differential, write $P_\ell(1)=L_\ell(E,1)^{-1}$,
$\tau_\ell=\length E(\Q_\ell)[p^\infty]$, and
$\log_{\omega_E}E(\Q_p)=p^l\Z_p$.  Then
\begin{equation}\label{models:haar-identities}
 \tau_\ell=v_p(c_\ell P_\ell(1))\quad(\ell\ne p),\qquad
 \tau_p-l=v_p(c_pP_p(1)).
\end{equation}
The connected real period $\Omega_0$ on a primitive positive Betti
cycle satisfies $\Omega_E/\Omega_0=\#\pi_0(E(\R))$, a $p$-adic unit.
\end{lemma}

\begin{proof}
The measure of the identity-component points for the minimal
differential is the nonsingular residue size divided by $\ell$,
namely $P_\ell(1)$.  Multiplication by the component number gives
$c_\ell P_\ell(1)$.  Away from $p$, the formal subgroup has pro-$\ell$
order and contributes no $p$-primary valuation.  At $p$, a sufficiently
small subgroup is identified by formal logarithm with its additive
lattice.  The $p$-primary part of the logarithm's finite kernel has
size $p^{\tau_p}$, and its image has measure $p^{-l}$; the
prime-to-$p$ kernel contributes a valuation unit.  This gives the
second identity.  The real
identity follows by translation of the minimal differential across
the components.
\end{proof}

\begin{lemma}[Global Kummer volume]\label{models:kummer-volume}
Let $K$ be imaginary quadratic, suppose $E(K)$ has rank one and
$\Sha(E/K)$ is finite, and let $C_{\rm Kum}$ impose the full integral
Kummer conditions.  Put $\tau_g=v_p(\#E(K)_{\rm tors})$ and
$s_K=v_p(\#\Sha(E/K))$.  For a nontorsion point $P$ let $n_P$ be
the valuation of its index in $E(K)/E(K)_{\rm tors}$.  Then
\begin{equation}\label{models:kummer-volume-formula}
 d(C_{\rm Kum},(P,P^\vee))=2n_P+2\tau_g-s_K.
\end{equation}
\end{lemma}

\begin{proof}
Derived Kummer reduction gives the usual finite Selmer groups;
passage to divisible coefficients gives the usual Kummer sequence.
Whole-Sha finiteness therefore identifies integral degree one with
the compact Mordell--Weil module.  The torsion in degree two is the
finite quotient in the compact-to-divisible comparison and has
length $s_K$.  Duality gives torsion length $\tau_g$ in degree three,
as well as a unimodular pairing on the free parts in degrees one and
two.  The torsion lengths in degrees one, two, three are consequently
$\tau_g,s_K,\tau_g$.  Their inverse-determinant contribution is
$2\tau_g-s_K$, and the point and its pairing functional each
contribute $n_P$.  This proves the formula, including under
conjugate-field duality; compare \cite[Section 3, ``Central Kummer
lattices and local measures'']{B}.
\end{proof}

\subsection{Bounded series and specialization}

A bounded series has coefficients in a fixed complete DVR, or in its
fraction field with one common denominator.  This condition excludes
general formal power series over the fraction field.  For measure
coefficients we also use
\[
 \mathcal V=W(\overline{\F}_p),\qquad
 \mathcal V^*=\varprojlim_j\prod_{\mathcal U}(\mathcal V/p^j)
 =W\left(\prod_{\mathcal U}\overline{\F}_p\right).
\]
This is again a complete integer DVR; fixed finite ramification may
be adjoined.  Arithmetic model matrices are first formed over the
smaller rings with finite residual coefficients.  Bounded measures
may then use this larger constant ring.  Fixed character substitution
and fixed finite group quotients commute with the limits.  A tame
cyclic generator is denoted by $1+u$; its relation disappears at every
fixed Artin precision when its $p$-power order tends to infinity.

\begin{lemma}[Division from high characters]\label{models:series-division}
Let $\mathcal O$ be one of these complete integer DVRs and let
$L,B\in\mathcal O[[t,\mathbf u]]$.  Suppose
$L\not\equiv0\pmod{(\varpi,t)}$, where $\varpi$ is a uniformizer.
If, for every sufficiently high finite $p$-power character $\theta$
of $1+t$, the specialization $B(\theta,\mathbf u)$ is divisible by
$L(\theta,\mathbf u)$ after inversion of $p$, then $L$ divides $B$
integrally.  The assertion applies simultaneously to finite products
of such divisors and after the fixed constant extensions above.
\end{lemma}

\begin{proof}
Choose integers $n_j$, one equal to $1$ and the others with suitably
large successive $p$-adic valuations, so that restriction to
$1+u_j=(1+x)^{n_j}$ separates a nonzero leading monomial of the
reduction.  Extend this to a group-power coordinate change.  Then
$L$ is distinguished, up to a unit, in one tame coordinate.
Applying this to a product handles a finite list.  Prescribed
characteristic-zero generic ranks and nonvanishings at $t=0$ may also
be retained: their leading homogeneous forms in logarithmic
coordinates impose finitely many proper conditions on the integer
choices, which can be avoided within the specified congruences.

Divide $B$ by the monic Weierstrass polynomial of $L$.  At each high
$\theta$ its distinguished degree is unchanged, and its quotient
module is $p$-torsion-free.  Thus divisibility after inversion of $p$
makes the integral remainder zero.  Every coefficient of this
remainder, considered as a bounded series in $t$, vanishes at all
high finite characters and is therefore zero by one-variable
Weierstrass preparation.  The argument also applies componentwise
to finite coefficient orders embedded in products of DVRs when the
same residual regularity is retained.  It is the proof of
\cite[Lemma 3.10]{B}, with $p$
in place of two.
\end{proof}

Two further forms of this principle will be useful.  First, a
nonzero bounded one-variable series has, at sufficiently high
characters, valuation equal to its minimum coefficient valuation
plus the first degree attaining that minimum times the valuation
of $\theta-1$.  This is unique dominance; it also proves finite
character avoidance for several variables by iteration.

\begin{lemma}[Persistence of poles]\label{models:poles}
Suppose integral stage numerators and denominators in a genuine
variable converge coefficientwise along the ultrafilter, with
nonzero limiting denominator.  If every stage numerator is
divisible by its denominator after inversion of $p$, each zero
of the limiting denominator in the open unit disk is a zero
of the limiting numerator with at least the same multiplicity.
\end{lemma}

\begin{proof}
Coefficientwise converging integral series converge uniformly on
smaller closed disks.  A zero of multiplicity $m$ of a nonzero
limiting denominator can be isolated on such a disk.  Weierstrass
preparation there retains its zero count on an ultrafilter-large
set.  If stage numerators are divisible by stage denominators after
inversion of $p$, the limiting numerator has at least that
multiplicity.  This pole test permits uncontrolled stage
denominators because boundedness is imposed on numerator and
denominator separately.  Tame characters, when present, are fixed
before this one-variable argument.
\end{proof}

\begin{lemma}[Points on a horizontal divisor]\label{models:field-points}
For a characteristic-zero height-one divisor in a bounded
$\Z_p$-power-series ring, finitely many functions nonzero on that
divisor can be kept nonzero at a point whose coordinates lie in the
open unit disk of a finite extension of $\Q_p$.
\end{lemma}

\begin{proof}
Its prime generator is not divisible by $p$.  A weighted coordinate
change makes it distinguished in one variable.  Modulo its
Weierstrass polynomial the norms of the finitely many nonzero
functions to the remaining power-series ring are nonzero.  Choose
the other coordinates, in finite extensions and in their open
unit disks, avoiding these norms; bounded-series avoidance gives
such a choice.  A root of the distinguished polynomial is also
interior and belongs to a finite extension.  It gives the required
point.
\end{proof}

\begin{remark}[Central determinant coordinates]\label{models:central-pivots}
When the generic and specialized cohomology dimensions agree, a
determinant coordinate specializes using any surviving set of
Gaussian pivots.  One need not keep a particular rational pivot
formula whose denominator vanishes at the point.  At actual central
stage calculations, bounds on torsion lengths and the corresponding
rank tests keep suitable minors nonzero in the limit.  Numerator,
denominator, and class coordinates then converge by their finite
matrix formulas.  This is why retaining torsion in the integral
models is necessary for central valuations.
\end{remark}
